% GENERATED FILE. Edit canonical lessons, then run npm run generate:books.
% canonical-source-hash: fnv1a64:1ba7d0dc66eb4b51
% canonical-lessons: GU-C502-karvat, GU-C502-korvu, GU-C502-hathodi, GU-C502-mekh, GU-C502-thokvu

\chapter{In the Carpenter's Shop}
\label{ch:gu-in-the-carpenter-s-shop}
\hlchaptermodality{502}

\begin{chapteropening}
\textbf{By the end of this chapter:} \emph{I can say how a carpenter saws wood and hammers in nails.}

Complete the last lesson of chapter 502: I can say how a carpenter saws wood and hammers in nails.
\end{chapteropening}

\section[\texorpdfstring{\gu{કરવત} (karvat)}{કરવત (karvat)}]{\gu{કરવત} (karvat) — a saw}

\label{lesson:GU-C502-karvat}



\begin{quote}
Before the new one: say the Gujarati for to know how, then the Gujarati for a profession.
\end{quote}

\subsection*{You'll want to know: \gu{કરવત}}
\textbf{\gu{કરવત}} — \emph{karvat} — \textquotedblleft{}a saw\textquotedblright{}.

\begin{quote}
\gu{સુથાર} \gu{કરવતથી} \gu{લાકડું} \gu{કાપે} \gu{છે}.
\end{quote}

\begin{quote}
\emph{suthār karvatthī lākḍũ kāpe che.}
\end{quote}

\begin{quote}
The carpenter cuts the wood with a saw.
\end{quote}

A neuter noun.

\begin{cousinweb}
\textbf{\gu{કરવત}} (\emph{karvat}) is inherited from Sanskrit \emph{karapattra}, \textquotedblleft{}a saw\textquotedblright{}, and CDIAL lists the Gujarati word (CDIAL 2795).
\end{cousinweb}

\begin{grammarlens}[title={What you've built}]
One more word for describing what you lived and felt.
\end{grammarlens}

\subsection*{Your turn}
\begin{itemize}
\raggedright
  \item \emph{Say it:} \emph{karvat}
  \item \emph{Say it:} \emph{karvat}, once more
  \item \emph{Say it:} \emph{suthār karvatthī lākḍũ kāpe che.} — The carpenter cuts the wood with a saw.
  \item \emph{Recall:} say the Gujarati for to know how, then the Gujarati for a profession, then say \emph{karvat} again
\end{itemize}

\begin{tcolorbox}[breakable,colback=teal!4,colframe=teal!35!black,title={Before you move on}]
What does \gu{કરવત} mean? (\textquotedblleft{}A saw\textquotedblright{}.) Say it once more.
\end{tcolorbox}
\section[\texorpdfstring{\gu{કોરવું} (korvũ)}{કોરવું (korvũ)}]{\gu{કોરવું} (korvũ) — to hollow out (wood)}

\label{lesson:GU-C502-korvu}



\begin{quote}
Before the new one: say the Gujarati for a profession, then the Gujarati for a saw.
\end{quote}

\subsection*{You'll want to know: \gu{કોરવું}}
\textbf{\gu{કોરવું}} — \emph{korvũ} — \textquotedblleft{}to hollow out (wood)\textquotedblright{}.

\begin{quote}
\gu{સુથારે} \gu{લાકડું} \gu{કોરીને} \gu{એક} \gu{વાટકી} \gu{બનાવી}.
\end{quote}

\begin{quote}
\emph{suthāre lākḍũ korīne ek vāṭkī banāvī.}
\end{quote}

\begin{quote}
The carpenter hollowed out the wood and made a bowl.
\end{quote}

Past, with \gu{તેણે} (he or she …): \gu{તેણે} \gu{લાકડું} \gu{કોર્યું}, \emph{teṇe lākḍũ koryũ}.

\begin{cousinweb}
\textbf{\gu{કોરવું}} (\emph{korvũ}) goes back to a reconstructed \emph{*kōrati}, \textquotedblleft{}scrapes\textquotedblright{}, and CDIAL lists the Gujarati verb for scooping out, carving or boring a hole (CDIAL 3530).
\end{cousinweb}

\begin{grammarlens}[title={What you've built}]
One more word for describing what you lived and felt.
\end{grammarlens}

\subsection*{Your turn}
\begin{itemize}
\raggedright
  \item \emph{Say it:} \emph{korvũ}
  \item \emph{Say it:} \emph{korvũ}, once more
  \item \emph{Say it:} \emph{suthāre lākḍũ korīne ek vāṭkī banāvī.} — The carpenter hollowed out the wood and made a bowl.
  \item \emph{Recall:} say the Gujarati for a profession, then the Gujarati for a saw, then say \emph{korvũ} again
\end{itemize}

\begin{tcolorbox}[breakable,colback=teal!4,colframe=teal!35!black,title={Before you move on}]
What does \gu{કોરવું} mean? (\textquotedblleft{}To hollow out (wood)\textquotedblright{}.) Say it once more.
\end{tcolorbox}
\section[\texorpdfstring{\gu{હથોડી} (hathoḍī)}{હથોડી (hathoḍī)}]{\gu{હથોડી} (hathoḍī) — a hammer}

\label{lesson:GU-C502-hathodi}



\begin{quote}
Before the new one: say the Gujarati for a saw, then the Gujarati for to hollow out.
\end{quote}

\subsection*{You'll want to know: \gu{હથોડી}}
\textbf{\gu{હથોડી}} — \emph{hathoḍī} — \textquotedblleft{}a hammer\textquotedblright{}.

\begin{quote}
\gu{સુથારની} \gu{હથોડી} \gu{ટેબલ} \gu{પર} \gu{છે}.
\end{quote}

\begin{quote}
\emph{suthārnī hathoḍī ṭebal par che.}
\end{quote}

\begin{quote}
The carpenter's hammer is on the table.
\end{quote}

A feminine noun.

\begin{cousinweb}
\textbf{\gu{હથોડી}} (\emph{hathoḍī}) comes from \emph{*hastakūṭa}, \textquotedblleft{}a hand hammer\textquotedblright{}, a form CDIAL builds from \emph{hásta}, \textquotedblleft{}hand\textquotedblright{}, and \emph{kūṭa}, \textquotedblleft{}a hammer\textquotedblright{}; CDIAL lists the Gujarati word (CDIAL 14028).
\end{cousinweb}

\begin{grammarlens}[title={What you've built}]
One more word for describing what you lived and felt.
\end{grammarlens}

\subsection*{Your turn}
\begin{itemize}
\raggedright
  \item \emph{Say it:} \emph{hathoḍī}
  \item \emph{Say it:} \emph{hathoḍī}, once more
  \item \emph{Say it:} \emph{suthārnī hathoḍī ṭebal par che.} — The carpenter's hammer is on the table.
  \item \emph{Recall:} say the Gujarati for a saw, then the Gujarati for to hollow out, then say \emph{hathoḍī} again
\end{itemize}

\begin{tcolorbox}[breakable,colback=teal!4,colframe=teal!35!black,title={Before you move on}]
What does \gu{હથોડી} mean? (\textquotedblleft{}A hammer\textquotedblright{}.) Say it once more.
\end{tcolorbox}
\section[\texorpdfstring{\gu{મેખ} (mekh)}{મેખ (mekh)}]{\gu{મેખ} (mekh) — a nail (for wood or a wall)}

\label{lesson:GU-C502-mekh}



\begin{quote}
Before the new one: say the Gujarati for to hollow out, then the Gujarati for a hammer.
\end{quote}

\subsection*{You'll want to know: \gu{મેખ}}
\textbf{\gu{મેખ}} — \emph{mekh} — \textquotedblleft{}a nail (for wood or a wall)\textquotedblright{}.

\begin{quote}
\gu{દીવાલમાં} \gu{એક} \gu{મેખ} \gu{છે}.
\end{quote}

\begin{quote}
\emph{dīvālmā̃ ek mekh che.}
\end{quote}

\begin{quote}
There is a nail in the wall.
\end{quote}

A feminine noun.

\begin{cousinweb}
\textbf{\gu{મેખ}} (\emph{mekh}) is usually traced to Persian \emph{mīkh}, \textquotedblleft{}a nail, a peg\textquotedblright{}.
\end{cousinweb}

\begin{grammarlens}[title={What you've built}]
One more word for describing what you lived and felt.
\end{grammarlens}

\subsection*{Your turn}
\begin{itemize}
\raggedright
  \item \emph{Say it:} \emph{mekh}
  \item \emph{Say it:} \emph{mekh}, once more
  \item \emph{Say it:} \emph{dīvālmā̃ ek mekh che.} — There is a nail in the wall.
  \item \emph{Recall:} say the Gujarati for to hollow out, then the Gujarati for a hammer, then say \emph{mekh} again
\end{itemize}

\begin{tcolorbox}[breakable,colback=teal!4,colframe=teal!35!black,title={Before you move on}]
What does \gu{મેખ} mean? (\textquotedblleft{}A nail (for wood or a wall)\textquotedblright{}.) Say it once more.
\end{tcolorbox}
\section[\texorpdfstring{\gu{ઠોકવું} (ṭhokvũ)}{ઠોકવું (ṭhokvũ)}]{\gu{ઠોકવું} (ṭhokvũ) — to hammer in}

\label{lesson:GU-C502-thokvu}



\begin{quote}
Before the new one: say the Gujarati for a hammer, then the Gujarati for a nail.
\end{quote}

\subsection*{You'll want to know: \gu{ઠોકવું}}
\textbf{\gu{ઠોકવું}} — \emph{ṭhokvũ} — \textquotedblleft{}to hammer in\textquotedblright{}.

\begin{quote}
\gu{અરુણે} \gu{હથોડીથી} \gu{દીવાલમાં} \gu{મેખ} \gu{ઠોકી}.
\end{quote}

\begin{quote}
\emph{Aruṇe hathoḍīthī dīvālmā̃ mekh ṭhokī.}
\end{quote}

\begin{quote}
Arun hammered a nail into the wall.
\end{quote}

Past, with \gu{તેણે} (he or she …): \gu{તેણે} \gu{મેખ} \gu{ઠોકી}, \emph{teṇe mekh ṭhokī}.

\begin{cousinweb}
\textbf{\gu{ઠોકવું}} (\emph{ṭhokvũ}) comes from \emph{*ṭhōkk}, \textquotedblleft{}to knock\textquotedblright{}, a form CDIAL reconstructs from its descendants, and CDIAL lists the Gujarati verb, \textquotedblleft{}to strike\textquotedblright{} (CDIAL 5513). The past agrees with what is hammered in: \textbf{\gu{મેખ} \gu{ઠોકી}}.
\end{cousinweb}

\begin{grammarlens}[title={What you've built}]
One more word for describing what you lived and felt.
\end{grammarlens}

\subsection*{Your turn}
\begin{itemize}
\raggedright
  \item \emph{Say it:} \emph{ṭhokvũ}
  \item \emph{Say it:} \emph{ṭhokvũ}, once more
  \item \emph{Say it:} \emph{Aruṇe hathoḍīthī dīvālmā̃ mekh ṭhokī.} — Arun hammered a nail into the wall.
  \item \emph{Recall:} say the Gujarati for a hammer, then the Gujarati for a nail, then say \emph{ṭhokvũ} again
\end{itemize}

\begin{tcolorbox}[breakable,colback=teal!4,colframe=teal!35!black,title={Before you move on}]
What does \gu{ઠોકવું} mean? (\textquotedblleft{}To hammer in\textquotedblright{}.) Say it once more.
\end{tcolorbox}
